% Source-derived chapter 04: The motive of an abelian variety
% Original source: standard-conjectures-abelian-fourfolds
\chapter{The motive of an abelian variety}
\label{standard-conjectures-abelian-fourfolds-chapter-04}
\source{standard-conjectures-abelian-fourfolds}

\begin{remark}[Source section]
\label{standard-conjectures-abelian-fourfolds-chapter-04-source}
\source{standard-conjectures-abelian-fourfolds}
\source{standard-conjectures-abelian-fourfolds}
This chapter reproduces the corresponding section of the retrieved
LaTeX source, including its definitions, statements, examples, and
proofs. It has not yet been translated into Lean.
\notready
\end{remark}

% The source section title was: The motive of an abelian variety
\label{reminder}
\source{standard-conjectures-abelian-fourfolds}


%Note that, following these references, the realization of a motive is a graded vector space.

%If not explicitly stated, we will work with general Weil cohomologies (not necessary classical ones).


In this section we recall classical results on motives of abelian type. We will work with the category of Chow motives $\CHM(k)_{F}$, we fix a Weil cohomology $H^*$ together with its realization functor $R$. Generalities on this category   can be found in the  Conventions. 
%The following theorem summarizes the results about motives of abelian varieties which will be used.
\begin{theorem}
\source{standard-conjectures-abelian-fourfolds}
\notready\label{classic} Let $A$ be an abelian variety of dimension  $g$. Let $\End(A)$ be its ring of endomorphisms (as group scheme) and $\mathfrak{h}(A)\in \CHM(k)_{F}$ be its motive. Then the following holds:
\source{standard-conjectures-abelian-fourfolds}
\begin{enumerate}
\item\label{kunneth} \cite{DeMu} The motive $\mathfrak{h}(A)$ admits a Chow--K\"unneth decomposition 
\source{standard-conjectures-abelian-fourfolds}
\[\mathfrak{h}(A)=\bigoplus_{n=0}^{2g} \mathfrak{h}^n(A)\] natural in $\End(A)$ and such that 
\[R(\mathfrak{h}^n(A))= H^n(A).\] 
\item\label{kunn} \cite{Ku1} The intersection product induces a canonical isomorphism of graded algebras \[\mathfrak{h}^*(A)=\Sym^* \mathfrak{h}^1(A).\]
\source{standard-conjectures-abelian-fourfolds}
\item\label{kings}   \cite[Proposition 2.2.1]{Ki}  The action of $\End(A)$ on $\mathfrak{h}^1(A)$ (coming from naturality in (\ref{kunneth}))  induces an injective morphism of algebras  \[\End(A)\otimes_{\Z}F \hookrightarrow \End_{\CHM(k)_{F}}(\mathfrak{h}^1(A))\]
\source{standard-conjectures-abelian-fourfolds}
and if $A$ is isogenous to $B\times C$ then $\mathfrak{h}^1(A) =  \mathfrak{h}^1(B) \oplus \mathfrak{h}^1(C) .$
\item\label{polarization}   \cite{Ku2} The Chow--Lefschetz conjecture holds true for $A$. In particular, the classical isomorphism in $\ell$-adic cohomology  induced by a polarization $H_{\ell}^1(A)\cong H_{\ell}^1(A)^{\vee}(-1)$ lifts to an isomorphism
\source{standard-conjectures-abelian-fourfolds}
 \[\mathfrak{h}^1(A)\cong\mathfrak{h}^1(A)^{\vee}(-1).\] 
\end{enumerate}
\end{theorem}
\begin{remark}
\source{standard-conjectures-abelian-fourfolds}
\notready\label{AHP}
\source{standard-conjectures-abelian-fourfolds}
We will need the above results also  in a slightly more general context, namely over the ring of integers of a $p$-adic field. Nowadays these results are known over very general bases, see \cite[Theorem 5.1.6]{OS2} or \cite[Theorem 3.3]{AHP}.
\end{remark}
\begin{definition}
\source{standard-conjectures-abelian-fourfolds}
\notready\label{defrank}
\source{standard-conjectures-abelian-fourfolds}
 A motive is called of abelian type if it  is a  direct factor of the motive of an abelian variety (up to Tate twist).
 
 
We say that a motive of abelian type $T$ is  of rank $d$  if the  cohomology groups of $R(T)$ are all zero except in one degree and in that degree the cohomology group is of dimension $d$.  In this case we will write $\dim T = d.$
\end{definition}

\begin{remark}
\source{standard-conjectures-abelian-fourfolds}
\notready\label{remarkrank}
\source{standard-conjectures-abelian-fourfolds}
For motives of abelian type this definition is known to be independent of $R$, see for example  \cite[Corollary 3.5]{jannsen} and \cite[Corollary 1.6]{Ascona}.
\end{remark}
%\begin{lemma}\label{lemmanumerical}
%Let us keep the same assumptions as Theorem \ref{mainthm} except the dimensional assumption in (\ref{assumptiondim2}). On the space $\Hom_{\CHM(k)_\Q}(\one, M_{\vert{k}})$ an element $v$ is numerically nonzero if and only if there exists another element $w$ such that $q \circ (v \otimes w) $ is non zero. In particular $q_Z$ defines a perfect pairing on $V_Z$.
%\end{lemma}
\begin{prop}
\source{standard-conjectures-abelian-fourfolds}
\notready\label{motivisupersingolari} 
\source{standard-conjectures-abelian-fourfolds}
Let $T$ be a motive of abelian type of dimension $d$. 
Consider   its space of algebraic classes modulo numerical equivalence \[V_Z = V_Z(T) =  \Hom_{\NUM(k)_\Q}(\one, T).\] Then the  inequality 
$\dim_\Q V_Z \leq d$ holds.

Moreover, if the equality $\dim_\Q V_Z = d$ holds then we have the following facts:

\begin{enumerate}
\item All  realizations of $T$ are spanned by algebraic classes.
\item Numerical equivalence on $\Hom_{\CHM(k)_\Q}(\one,T)$ coincides with  homological equivalence (for all  cohomologies).
\item Call $L$ the field of coefficients of the realization $R$, then the equality $V_Z\otimes _\Q L = R(T)$ holds.
\end{enumerate}
%\begin{enumerate}
%\item[(a)]\label{massimicicli} $  \dim_\Q V_Z = d.$
%\item[(b)]  All  realizations of $X$ are spanned by algebraic classes.
%\item[(c)]  For one   realization $R$, the vector space $R(X)$ is spanned by algebraic classes.
%\end{enumerate}
%Moreover, under these equivalent conditions, the numerical equivalence on $\Hom_{\CHM(k)_\Q}(\one,X)$ coincides with the homological equivalence for all  cohomologies.
\end{prop}
\begin{proof}
Let us consider $n$ elements $\bar{v}_1,\ldots,\bar{v}_n$ which are linearly independent in $V_Z$ and fix $v_1,\ldots,v_n\in \Hom_{\CHM(k)_\Q}(\one,T)$ which are liftings of those. By definition of numerical equivalence, there exist  $n$ elements   $f_1,\ldots,f_n $  in $\Hom_{\CHM(k)_\Q}( T,\one)$ such that $f_i(v_j)=\delta_{ij}$. By applying a realization we get $R(f_i)(R(v_j))=\delta_{ij}$ which implies the inequality ${n \leq \dim R(T)=d.}$


When $n=d$, the argument just above shows that $R(v_1),\ldots,R(v_d)$ form a basis of $\Hom(R(\one),R(T))$, which means that $R(T)$ is spanned by the  algebraic classes $R(v_1)(1),\ldots,R(v_d)(1)$, hence we have (1).

Fix  an element $v\in \Hom_{\CHM(k)_\Q}(\one,T)$ and write its realization in the previous basis $R(v)= \lambda_1 R(v_1)+ \ldots + \lambda_d R(v_d)$. Consider now the composition $f_i(v)$. On the one hand, it is a rational number, on the other hand it equals $\lambda_i$. This means that any algebraic class is a rational combination of the basis $R(v_1)(1),\ldots,R(v_d)(1)$. This implies the points (2) and (3).
\end{proof}
